\documentclass[../DoAn.tex]{subfiles}
\begin{document}

\begin{center}
\Large{\textbf{ABSTRACT}}
\end{center}

\vspace{0.8cm}

This thesis develops a \textbf{citation-grounded question-answering system for Vietnamese labor law} based on Microsoft \gls{GraphRAG}~\cite{graphrag2024}, addressing the limitations of traditional keyword-based legal search~\cite{robertson1994bm25} (inability to perform cross-document reasoning) and general-purpose \gls{LLM} chatbots (prone to hallucination~\cite{ji2023hallucination} and outdated or incorrect citations).

The \textbf{main contribution} is a \textbf{two-layer ontology with canonicalization}: Layer~1 (deterministic) represents the formal legal structure without using an LLM, ensuring fully consistent Article/Clause citations; Layer~2 (LLM extraction) captures domain semantics with greater flexibility. Building on this design, the system extends to multi-hop reasoning (traversing legal relationships rather than relying on vector similarity), filtering of expired citations, and contract compliance analysis that combines deterministic rules with graph grounding.

Experimental results show that the proposed approach outperforms standalone LLMs and traditional vector-based \gls{RAG} on questions requiring accurate citations and cross-provision reasoning. The system achieves 100\% reasoning accuracy on the multi-hop question group (compared with 33\% for Local Search alone), demonstrating that deterministic legal structure combined with a two-layer ontology is a necessary foundation for reliable legal question answering.

\textbf{Limitations}: a narrow corpus (labor law only) and a small evaluation set not yet validated by legal experts. Nevertheless, the design pattern can be extended to other legal domains and reused in similar Legal AI applications. The system is intended solely as a decision-support tool and does not replace official legal consultation.

\vspace{1.5cm}

\begin{flushright}
\textbf{Author}\\[0.1cm]
\textit{(Signature and Full Name)}\\[2.5cm]
\textbf{Nguyen Dat Trong}
\end{flushright}

\end{document}
